% ============ 中文排版（XeLaTeX，无 xeCJK 的原生方案）============
% XeTeX 内建 CJK 断行（环境无 xeCJK）
\XeTeXlinebreaklocale "zh"
\XeTeXlinebreakskip = 0pt plus 0.08em minus 0.04em
\linespread{1.28}
\setlength{\parskip}{0.42em}
\setlength{\parindent}{0pt}
\frenchspacing
\sloppy
\hyphenpenalty=5000 \tolerance=2000 \emergencystretch=3em

% ============ 配色（与 HTML 版一致）============
\usepackage{xcolor}
\definecolor{bookcrimson}{HTML}{A51C30}
\definecolor{boxrule}{HTML}{D0D4D9}
\definecolor{boxback}{HTML}{FAFAFB}
\definecolor{fpback}{HTML}{FDF4F6}
\definecolor{fprule}{HTML}{A51C30}

% ============ 自定义盒子 ============
\usepackage[most]{tcolorbox}
\newtcolorbox{fpbox}{
  enhanced, breakable, colback=fpback, colframe=fprule,
  boxrule=0pt, leftrule=2.4pt, arc=1.6pt,
  left=10pt, right=10pt, top=7pt, bottom=7pt, before skip=10pt, after skip=10pt }
\newtcolorbox{purposebox}{
  enhanced, breakable, colback=white, colframe=white,
  boxrule=0pt, leftrule=2.4pt, colframe=bookcrimson,
  left=10pt, right=10pt, top=3pt, bottom=3pt, before skip=8pt, after skip=12pt,
  fontupper=\itshape\color{bookcrimson} }

% ============ 标题与版式 ============
\usepackage{titlesec}
\titleformat{\chapter}[display]
  {\sffamily\bfseries\color{bookcrimson}}
  {\fontsize{48}{48}\selectfont\thechapter}
  {8pt}{\fontsize{22}{26}\selectfont\raggedright}
\titlespacing*{\chapter}{0pt}{-24pt}{28pt}
\titleformat*{\section}{\sffamily\bfseries\large\color{bookcrimson}}
\titleformat*{\subsection}{\sffamily\bfseries\normalsize}
\titleformat*{\subsubsection}{\sffamily\bfseries\small}

\usepackage{fancyhdr}
\pagestyle{fancy}\fancyhf{}
\fancyhead[LE]{\small\sffamily\nouppercase{\leftmark}}
\fancyhead[RO]{\small\sffamily\nouppercase{\rightmark}}
\fancyfoot[C]{\small\sffamily\thepage}
\renewcommand{\headrulewidth}{0.3pt}

% 图表题注
\usepackage{caption}
\captionsetup{font=small, labelfont={bf,sf,color=bookcrimson}, labelsep=period,
              justification=raggedright, singlelinecheck=false, skip=6pt}

% 索引条目：无 makeidx 时让 \index 安全失效
\renewcommand{\index}[1]{}

% 中文无真斜体：\emph 改用无衬线体区分（替代西文斜体的惯例做法）
\let\origemph\emph
\renewcommand{\emph}[1]{{\sffamily #1}}
\renewcommand{\textit}[1]{{\sffamily #1}}

% microtype：字距与边缘对齐微调（同 MLSysBook）
\usepackage{microtype}

% ---------------------------------------------------------------- 整页封面
% build-pdf.sh 会把 cover-a4.jpg 拷到 .tex 同目录；文件不存在时静默跳过，
% 这样单独跑 quarto render --to pdf 也不会报错。
\usepackage{eso-pic}
\usepackage{graphicx}
\makeatletter
\IfFileExists{cover-a4.jpg}{%
  \let\bk@origmaketitle\maketitle
  \renewcommand{\maketitle}{%
    \begingroup
      \thispagestyle{empty}%
      \AddToShipoutPictureBG*{%
        \put(0,0){\includegraphics[width=\paperwidth,height=\paperheight]{cover-a4.jpg}}}%
      \null\clearpage
    \endgroup
    \bk@origmaketitle}%
}{}
\makeatother

% 拉丁扩展字符（参考文献人名等）回退字体，避免缺字
\usepackage{newunicodechar}
\newfontfamily\latinext{DejaVu Serif}
\newunicodechar{Ć}{{\latinext Ć}}
\newunicodechar{ć}{{\latinext ć}}
